\chapter{Exact Galerkin Computation}
\label{chap:exact-galerkin-computation}

The computation closes when the weak form becomes a report. The report should show the state, the raw residual,
the Sobolev residual, the squared Sobolev norm, and the values of the first Fréchet variation on selected tests.
If the residual is exactly zero, the theorem from the previous chapter supplies the weak least-activity
certificate. If it is not zero, the report says where the failure lives.

\section{The Toy Tensor}

The companion file includes a two-row exact toy system:

\begin{quote}\small
\begin{verbatim}
def diracToyTensor : UniversalTensor :=
  { diracRows :=
      [[1, -1],
       [1,  1]]
    tensorRows :=
      [[0, 0],
       [0, 0]]
    sobolevRows :=
      [[1, 0],
       [0, 1]]
    load :=
      [0, 2] }

def exactState : Vec :=
  [1, 1]
\end{verbatim}
\end{quote}

The first row reads \(c_0-c_1\). The second row reads \(c_0+c_1\). The load is \((0,2)\), so the state
\((1,1)\) closes both rows exactly. The Sobolev reader is the identity in this toy example. A real run would
replace the rows with assembled Dirac, tensor, boundary, and Sobolev matrices. The shape of the computation
would not change.

\section{The Report}

The report object is:

\begin{quote}\small
\begin{verbatim}
structure GalerkinReport where
  state : Vec
  rawResidual : Vec
  sobolevResidual : Vec
  sobolevNormSq : Int
  frechetOnTests : List Int
deriving Repr
\end{verbatim}
\end{quote}

The report is produced by:

\begin{quote}\small
\begin{verbatim}
def report
    (U : UniversalTensor) (state : Vec) (tests : List Vec) :
    GalerkinReport :=
  { state := state
    rawResidual := rawTensorResidual U state
    sobolevResidual := sobolevResidual U state
    sobolevNormSq := sobolevNormSq U state
    frechetOnTests := tests.map (fun test =>
      firstFrechetVariation U state test) }
\end{verbatim}
\end{quote}

Run it from the Lean project:

\begin{quote}\small
\begin{verbatim}
cd device
lake env lean Measurement/WeakDiracGalerkin.lean
\end{verbatim}
\end{quote}

The example prints:

\begin{quote}\small
\begin{verbatim}
{ state := [1, 1],
  rawResidual := [0, 0],
  sobolevResidual := [0, 0],
  sobolevNormSq := 0,
  frechetOnTests := [0, 0, 0] }
\end{verbatim}
\end{quote}

That is the finite weak solution in its smallest visible form. The raw residual closes. The Sobolev residual
closes. The Sobolev norm square is zero. The first Fréchet variation vanishes on the listed tests, and the
theorem says exact residual closure makes it vanish on every finite test vector in the skeleton.

\section{Replacing the Toy Rows}

For a serious Dirac/Galerkin experiment, the user replaces the toy rows in four stages.

\begin{enumerate}
\item Choose the finite trial and test basis: spinor components, element shape functions, boundary functions,
and any gauge degrees of freedom.
\item Assemble the Dirac rows from \(i\gamma^\mu(\partial_\mu+iqA_\mu)-m\), including boundary and connection
data at the chosen floor.
\item Assemble the universal tensor rows from the zeroth/tensor term: metric, frame, curvature, boundary
anchor, or other tensor carrier required by the experiment.
\item Assemble the Sobolev reader: mass matrix, stiffness matrix, mass-plus-stiffness matrix, or another exact
row system representing the chosen \(H^s\) norm.
\item Enter the load vector \(f\), compute the residual, and ask Lean for the exact report.
\end{enumerate}

The finite method is exact only after these choices are fixed. Changing the basis changes the rows. Changing the
Sobolev reader changes the activity. Changing the boundary condition changes the weak problem. This is not a
weakness of Galerkin computation. It is the point: the computation says precisely which finite experiment was
performed.

\section{What Has Been Demonstrated}

The addition demonstrates the mathematical spine the experimentation volume needs:
\[
  \boxed{
  \hbox{weak least activity}
  =
  D\mathcal A_U(\psi)[\eta]=0
  =
  \langle R_U(\psi),\eta\rangle_{H^s}=0
  }
\]
with
\[
  R_U(\psi)=\mathcal D_A\psi+\mathcal T_U\psi-f .
\]
The Dirac operator supplies the first-order Clifford/gauge term. The universal tensor supplies the zeroth
background term. The Sobolev norm supplies the reader. The Fréchet derivative supplies the weak variation. The
Galerkin basis supplies the finite exact computation.

This is serious mathematics because it says where each burden lives. The continuum problem lives in the choice
of Sobolev space, domain, boundary, and convergence theorem. The finite experiment lives in the assembled rows.
The Lean theorem lives at the seam: exact finite Galerkin closure implies the weak first-variation equation for
the finite tests. The apparatus does not borrow a miracle from the continuum. It computes the rows it has, proves
what those rows imply, and leaves the stronger analytic claims named as the next obligations.

\section{The Final Floor}

The last floor is therefore not that Lean has solved the Dirac equation on a manifold. It has not. The last floor
is sharper and more useful: once the Dirac/tensor residual has been expressed in a Sobolev reader and assembled
on a finite Galerkin basis, the weak principle of least activity is an exact finite statement. The apparatus can
compute that statement without approximation. Approximation enters only when the finite result is read back
toward the continuum.

That is the handoff from experimentation to analysis. The experiment writes the rows. Lean checks the weak
closure. Analysis then asks how the finite closures behave under refinement. Each job has its own floor, and the
floors no longer blur.
